%-*-latex-*-
%-*-latex-*-
\newcommand\COURSE{ciss240}
\newcommand\ASSESSMENT{q2101}
\newcommand\ASSESSMENTTYPE{Quiz}
\newcommand\POINTS{\textwhite{xxx/xxx}}

\input{myquizpreamble}
\input{yliow}
\input{\COURSE}
%-*-latex-*-

%-*-latex-*-

%-*-latex-*-
\renewcommand\TITLE{\ASSESSMENTTYPE \ \ASSESSMENT}

\renewcommand\EMAIL{}
\renewcommand\AUTHOR{eagraham1@cougars.ccis.edu}

\textwidth=6in
\begin{document}
\topmatterthree

This is a closed-book, no compiler, 5 minute quiz.

The following questions involve writing function prototypes.
\textsc{Warning:} Use \verb!&! only when absolutely necessary.

%------------------------------------------------------------------------------
\nextq
The function \verb!swap! swaps the values of two \verb!int! variables.
For instance if \verb!x! is \verb!0! and \verb!y! is \verb!1!,
on returning from \verb!swap(x, y)!, \verb!x! is \verb!1! and
\verb!y! is \verb!0!.
Write down the function prototype of \verb!swap!.
Do not include variable names in the function prototype.
\\
\ANSWER
\begin{answercode}
void swap(int&, int&);
\end{answercode}

%------------------------------------------------------------------------------
\nextq
The function \verb!swap_double! swaps the values of two
\verb!double! variables.
For instance is \verb!x! is \verb!0.0! and \verb!y! is \verb!1.1!,
on returning from \verb!swap_double(x, y)!, \verb!x! is \verb!1.1! and
\verb!y! is \verb!0.0!.
Write down the function prototype of \verb!swap_double!.
Do not include variable names in the function prototype.
\\
\ANSWER
\begin{answercode}
void swap_double(double&, double&);
\end{answercode}

%------------------------------------------------------------------------------
\nextq
Let \verb!x! and \verb!y! be integer variables.
On calling \verb!f(x, y)!, the average (as a \verb!double! value)
of \verb!x! and \verb!y! is returned, \verb!x! is set to \verb!0!,
\verb!y! is unchanged.
Write down the function prototype of \verb!f!.
Do not include variable names in the function prototype.
\\
\ANSWER
\begin{answercode}
double f(int&, int);
\end{answercode}

%------------------------------------------------------------------------------
\nextq
If \verb!x! is an integer array, \verb!n! is an integer variable,
\verb!value! is an integer variable,
then on calling \verb!fill(x, n, value)!,
the function will set
\verb!x[0]! to \verb!value!,
\verb!x[1]! to \verb!value!,
\verb!x[2]! to \verb!value!, ...,
\verb!x[n - 1]! to \verb!value!.
Write down the function prototype of \verb!fill!.
You must include variable names in the function prototype
(use \verb!x!, \verb!n!, \verb!value! for their names).
\\
\ANSWER
\begin{answercode}
void fill(int x[], int n, int value);
\end{answercode}

%------------------------------------------------------------------------------
\nextq
The function \verb!play_audio()! plays an audio file.
The parameters are:
\begin{enumerate}[nosep]
\li \verb!filename!: a C-string that is the name of the audio file to be loaded.
\li \verb!track_number!: an integer.
$0$ is the first track, $1$ is the second track, etc.
$-1$ is \lq\lq play all tracks".
\li \verb!repeat!: a boolean.
Track(s) will be played in a loop if \verb!repeat!
is \verb!true!.
\end{enumerate}
The function returns a $0$ if there are no errors,
a $-1$ if the file cannot be found using the \verb!filename!,
$-2$ if the \verb!track_number! is not $-1$ and the track number is not found
in the audio file.
Write down the prototype of this function.
You must include variable names in the function prototype
(use \verb!filename!, \verb!track_number!, \verb!repeat! for their names).
\\
\ANSWER
\begin{answercode}
int play_audio(const char* filename, int track_number, bool repeat);
\end{answercode}

%------------------------------------------------------------------------------
%-*-latex-*-
\newpage
\input{instructions}
\end{document}

